%% This is file `elsarticle-template-1-num.tex',
%%
%% Copyright 2009 Elsevier Ltd
%%
%% This file is part of the 'Elsarticle Bundle'.
%% ---------------------------------------------
%%
%% It may be distributed under the conditions of the LaTeX Project Public
%% License, either version 1.2 of this license or (at your option) any
%% later version.  The latest version of this license is in
%%    http://www.latex-project.org/lppl.txt
%% and version 1.2 or later is part of all distributions of LaTeX
%% version 1999/12/01 or later.
%%
%% Template article for Elsevier's document class `elsarticle'
%% with numbered style bibliographic references
%%
%% $Id: elsarticle-template-1-num.tex 149 2009-10-08 05:01:15Z rishi $
%% $URL: http://lenova.river-valley.com/svn/elsbst/trunk/elsarticle-template-1-num.tex $
%%
\documentclass[preprint, authoryear,11pt, 3p]{elsarticle}

%% Use the option review to obtain double line spacing
%% \documentclass[preprint,review,12pt]{elsarticle}

%% Use the options 1p,twocolumn; 3p; 3p,twocolumn; 5p; or 5p,twocolumn
%% for a journal layout:
%% \documentclass[final,1p,times]{elsarticle}
%% \documentclass[final,1p,times,twocolumn]{elsarticle}
%% \documentclass[final,3p,times]{elsarticle}
%% \documentclass[final,3p,times,twocolumn]{elsarticle}
%% \documentclass[final,5p,times]{elsarticle}
%% \documentclass[final,5p,times,twocolumn]{elsarticle}

%% The graphicx package provides the includegraphics command.
\usepackage{graphicx}
\usepackage{subfig}
\usepackage{caption}
%% The amssymb package provides various useful mathematical symbols
\usepackage{amssymb}
%% The amsthm package provides extended theorem environments
%% \usepackage{amsthm}

%% The lineno packages adds line numbers. Start line numbering with
%% \begin{linenumbers}, end it with \end{linenumbers}. Or switch it on
%% for the whole article with \linenumbers after \end{frontmatter}.
\usepackage{lineno}

%% natbib.sty is loaded by default. However, natbib options can be
%% provided with \biboptions{...} command. Following options arehttps://v2.overleaf.com/project/5ae9a183e6e381104e1a8531
%% valid:

%%   round  -  round parentheses are used (default)
%%   square -  square brackets are used   [option]
%%   curly  -  curly braces are used      {option}
%%   angle  -  angle brackets are used    <option>
%%   semicolon  -  multiple citations separated by semi-colon
%%   colon  - same as semicolon, an earlier confusion
%%   comma  -  separated by comma
%%   numbers-  selects numerical citations
%%   super  -  numerical citations as superscripts
%%   sort   -  sorts multiple citations according to order in ref. list
%%   sort&compress   -  like sort, but also compresses numerical citations
%%   compress - compresses without sorting
%%
%% \biboptions{comma,round}

% \biboptions{}

\journal{Socius: Sociological Research for a Dynamic World}

\begin{document}

\begin{frontmatter}

%% Title, authors and addresses

\title{Durable Taste Schemas and Patterns of Cultural Choice}

%% use the tnoteref command within \title for footnotes;
%% use the tnotetext command for the associated footnote;
%% use the fnref command within \author or \address for footnotes;
%% use the fntext command for the associated footnote;
%% use the corref command within \author for corresponding author footnotes;
%% use the cortext command for the associated footnote;
%% use the ead command for the email address,
%% and the form \ead[url] for the home page:
%%
%% \title{Title\tnoteref{label1}}
%% \tnotetext[label1]{}
%% \author{Name\corref{cor1}\fnref{label2}}
%% \ead{email address}
%% \ead[url]{home page}
%% \fntext[label2]{}
%% \cortext[cor1]{}
%% \address{Address\fnref{label3}}
%% \fntext[label3]{}


%% use optional labels to link authors explicitly to addresses:
%% \author[label1,label2]{<author name>}
%% \address[label1]{<address>}
%% \address[label2]{<address>}

%\author[label1]{Omar Lizardo}

%\address[label1]{Department of Sociology, University of California, Los Angeles, 375 Portola Plaza, 264 Haines Hall, Los Angeles, CA, 90095. Email: olizardo@soc.ucla.edu.}

\begin{abstract}
%% Text of abstract
Whether cultural shemas are malleable or durable is a key, but still unresolved, theoretical debate in cultural analysis. In this paper, I contribute to a recent line of work seeking to move this issue to the empirical terrain. To do so, I extract taste schemas from survey data on a representative sample of Americans collected in 2012 containing items on likes and dislikes of musical genres similar to those included in the 1993 General Social Survey culture module. Consistent with the durability position, I find remarkably similar taste schemas in 2012 as those uncovered in the 1993 General Social Survey data in previous work using comparable techniques. Going beyond previous work, I validate these schemas using independent behavioral evidence on musical listening choices and propose a novel way to think about the link between taste schemas and culture consumption behavior that can be generalized to other domains
\end{abstract}

\begin{keyword}
Omnivorousness  \sep Taste \sep Schemas \sep Cultural Choices \sep Univores
%% keywords here, in the form: keyword \sep keyword
%% MSC codes here, in the form: \MSC code \sep code
%% or \MSC[2008] code \sep code (2000 is the default)
\end{keyword}

\end{frontmatter}
\newpage
\section{Introduction}
\subsection{The Durability versus Malleability Debate in Cultural Analysis}
Are cultural schemas durable or malleable? This is a question on which cultural analysts drawing from different theoretical traditions split into different camps \citep{vaisey2016cultural}. One the one hand, we have analysts who talk about schemas as ``deeply internalized'' and ``durable'' cultural elements, and thus as having the potential to show continuity in relation to changes in the external environment \citep{bourdieu1990logic, strauss1997cognitive, patterson2014making}. On the other hand, we have analysts who conceptualize schemas as highly responsive to changes in the external institutional environment and processes, and thus as ``shifting,'' ``multiform'' \citep{dimaggio1997culture, swidler2005anchors, sewell2004concept}. 

Pitched as a theoretical discussion, this debate runs the danger of becoming a purely scholastic issue. Here I break with this approach and instead follow the call recently laid out by \citet{vaisey2016cultural} to move the issue into a more empirical terrain. I use the case of cultural taste schemas to examine the issue of whether we can observe continuity or transformation at the population level over a twenty-year period. The results reveal somewhat remarkable levels of schematic continuity in this domain, a pattern consistent with the relative durability thesis \citep{vaisey2016cultural}. However, partially consistent with the malleability position, the differences that we do observe can be readily traced to institutional transformations on the consumer and production side of musical genres, such as the aging of the Rock and Roll generation and the emergence of hip-hop and rap as a prototypical ``pop'' or ``industry'' genre. 
\subsection{The Case of Cultural Taste Schemas}

Cultural taste schemas represent a strategic research site with which to address the issue of schematic durability. This is for two reasons. First, both theoretical and methodological developments with regard to extraction of schemas from survey data have used this case, and the associated data source (the 1993 General Social Survey), as their primary testbed \citep{goldberg2011mapping, boutyline2017improving}. This means that we have a set of relatively well-established results documenting population-level variability on taste schemas in the American public at this time characterizing the nature of each evaluative schema, thus facilitating comparison. 

Second, the domain of musical taste is different from other institutional settings such as politics. In the political domain, we might expect continuity and durability in people's conceptions to be fixed by the relative durability of external institutional arrangements (e.g. the two-party system in the U.S.). And in fact, this is precisely what we observe. In a landmark study, \citet{baldassarri2014neither} find remarkable continuity in the ways Americans organize their understanding of political issues at the population level across twenty years (although they could not engage in apples to apples comparisons of the same items as we do here). A durability theorist might see this as strong evidence for their position, but a malleability theorist could point to continuity (even with change at the margins) in external institutional arrangements as a competing explanation. Regarding politics, therefore, continuity seems over-determined by external and internal mechanisms, making it hard to adjudicate between the durability and malleability positions. 

Musical taste, by way of contrast, represents a case in which radical transformations have occurred in terms of production, distribution, modes of consumption, as well as definitions of genre popularity and ``aesthetic mobility'' \citep{baumann2001intellectualization} since the early 1990s. Thus, if we were to expect, consistent with the malleability thesis, that schemas respond to radical structural and external changes in a given domain, then the domain of musical evaluation would be an ideal case in which to observe such transformations. On the other hand, observing relative continuity in this domain would support a strong version of the durability thesis \citep{vaisey2016cultural}. 

\subsection{Terminology and Approach}
It is important to get clear on the use of certain specialized terms that will recur through the manuscript. In what follows, I refer to the underlying schema responsible for a given pattern of responses to evaluative questions in a cultural taste survey as a \emph{taste schema}. Persons who share the same (or similar) taste schemas are members of a \emph{schematic class}. Persons in the same schematic class may provide seemingly different (or even opposed) responses to survey questions on cultural taste. However, these responses, even if different at the level of content, share the same understanding of the overall evaluative organization of the field of objects in question \citep{goldberg2011mapping}. As such, traditional methods that cluster respondents on overt evaluative responses (e.g. factor or latent class analysis) are not useful for eliciting underlying taste schemas.

The analyst must then settle on a method to extract schemas from survey responses. Following \citep{boutyline2017improving}, I refer to \emph{schematic class analysis}, as a set of methodological procedures designed to elicit schemas from ordered responses to a given set of objects in survey data (such as likes and dislikes of musical genres in the case of cultural taste). This is done by assigning respondents to disjoint schematic classes based on a measure of similarity that is sensitive to the fact that two persons may belong to the same schematic class even though their vector of responses are (usually monotonic) transformations of one another (and thus overtly different). The outcome of schematic class analysis is thus to group respondents based on sharing similar schemas. 

Not all methods of extracting schemas from survey responses are equally valid. {\em Correlational Class Analysis} is a variant of schematic class analysis that uses the absolute value of the Pearson correlation coefficient between the response vectors of different respondents as the relevant similarity criterion and, like relational class analysis, partitions the resulting (weighted) network of respondents, where links are based on the absolute value of the correlation, using a modularity maximization algorithm. \citet{boutyline2017improving} shows that the correlation approach outperforms other ways of computing response pattern similarity, and thus eliciting schematic class membership, such as Goldberg's \citeyear{goldberg2011mapping} use of the ``relationality'' measure in relational class analysis. Given this, I use correlational class analysis in what follows.  

\subsection{Previous Work Using Schematic Class Analysis}
Schematic class analysis, whether based on relationality or the correlation distance, has been used in a variety of applications to extract schematic classes from mass survey responses. These include political belief logics \citep{baldassarri2014neither, daenekindt2017people}, patterns of thought about the markets and the economy \citep{dimaggio2018searching, metzmcDonnell2018multiple}, and taste schemas about the link between science and religion \citep{dimaggio2017culture}. In addition, recent work by \citet{daenekindt2017structure} analyzes the structure of aesthetic dispositions among a sample of museum-goers, a topic close to that dealt with here. In this respect, schematic class analysis has already made major contributions to the social science literature across a wide range of topics. 

\subsection{Contributions of the Present Work}
Previous work using schematic class analysis suffers from validation limitations. In particular, previous studies are forced to rely on inspecting differences in the correlations between items within classes (internal validation) or the correlation between schematic class membership and demographic characteristics to both validate schematic class assignment, and make inferences about the possible effect of schemas on behavior. \citet{goldberg2011mapping} used internal, and demographic validation strategy in the original study using 1993 General Social Survey data.\citet{baldassarri2014neither} used a combination of internal, demographic, and criterion-based strategies (looking at how predictors differ in predicting partisanship across schematic classes). Similarly \citet{dimaggio2018searching} and \citet{dimaggio2017culture} use a combination of internal, demographic, and substantive validation strategies.\footnote{Note that this discussion is not related to validation of the \emph{number} of classes extracted by schematic class analysis methods such as RCA or CCA. In this respect, \citet{boutyline2017improving} has used multiple group SEM, and \citet{dimaggio2017culture} and \citet{dimaggio2018searching} have used network randomization methods coupled with cluster validation (e.g., gap statistic) techniques.}

This study adds to and extends previous work by using both the standard internal validation approach (based on within-class inter-item correlations) and also a measure of choice behavior that is directly related to the domain in question, namely, music listening habits. In this respect, my strategy comes closest to the that pursued by \citet{baldassarri2014neither}. This allows me to determine empirically the behavioral potential of schematic classes in a way that establishes external (predictive) validity within the musical taste domain. I incorporate insights from recent work in the analysis of item correlation networks \citep{epskamp2012qgraph}, and link recently developed formal measures of item position in the taste schema with behavioral predictive criteria. Promising results obtained with this approach recommend it as a useful strategy for future schematic class analysis work. 

\section{Data and Variables}
In the summer of 2012, I fielded a survey that was (partially) designed to replicate the General Social Survey 1993 Culture Module. This last is now a canonical dataset, providing the empirical basis for large and variegated set of analyses (and re-analyses) in the sociology of taste \citep[][among others]{boutyline2017improving, bryson1996anything,bryson1997univores, goldberg2011mapping, han2003unraveling, okada2017structure, tampubolon2008revisiting}. It can be said without exaggeration that the bulk of our knowledge of the cultural taste patterns of the American population, both in terms of schematic form and actual content, are due to this data source. As such, the addition of a more contemporary source on the musical taste patterns of the American population constitutes a valuable addition on its own right \citep{lizardo2016cultural, lizardo2016end, lizardo2015musical}.

The data used in this paper were collected by Survey Sample International (SSI), a private firm that specializes in sampling, data collection, and analysis. SSI managed recruitment and participation invitation tasks to generate a panel of adults from which our working sample was drawn. Survey respondents were selected from the panel for participation on the basis of age, gender, race, education, and geographic region to approximate a sample representative of the U.S. population, $N = 2,250$.

Like General Social Survey 1993, SSI 2012 included items assessing respondents’ likes and dislikes (as well as a middle category of ``mixed feelings'') for 20 categories of musical style: (1) classical or symphony and chamber, (2) opera or operetta, (3) jazz, (4) Broadway musicals or show tunes, (5) mood or easy listening, (6) big band or swing, (7) classic rock or oldies, (8) country, (9) bluegrass, (10) folk, (11) hymns or gospel, (12) Latin or Spanish/salsa, (13) rap or hip hop, (14) blues or R and B, (15) reggae, (16) top 40 or pop music, (17) contemporary rock, (18) indie or alternative rock, (19) dance or club/electronic, and (20) hard rock or heavy metal. These responses are the basis of the schematic class analysis. 

The survey also asked respondents to report in a separate item whether they regularly {\em listen} to each one of the genres listed. Responses to these items allow us to examine behavioral expression (and variation) across schematic classes separately from the items used to elicit schematic class membership.

\section{Correlational Class Analysis}
I follow the same data-analytic procedure as \citet{boutyline2017improving}, with two modifications. The first is relatively minor; given the nature of the taste evaluation items as limited ordinal response items, I use the polychoric correlation coefficient as a ``drop-in replacement'' \citep[386]{boutyline2017improving} for the Pearson correlation.\footnote{Polychoric correlations were computed using the function \texttt{cor\_auto} of version 1.5 of the \texttt{R} package \texttt{qgraph} \citep{epskamp2012qgraph}.}

The second is less trivial and has to do with the way I treated the ``not familiar with this genre" responses to genre evaluation questions.  \citet{goldberg2011mapping} recoded these to the ``mixed feelings'' (middle) response for all respondents who offered this response for six genres or less (dropping those who were not familiar with seven or more genres). \citet{boutyline2017improving} followed this same coding strategy for purposes of strict replication. 

Since ``not familiar" responses are a type of missing data, I follow recently recommended best practice and use multiple imputation rather than fixed value recoding. This strategy is more efficient and does not require us to choose an arbitrary cutoff to decided how many ``not familiar" responses are sufficient to drop a respondent from the analysis. The variables used in the multiple imputation procedure consisted of the other genres for which the respondent did provide a valid taste response. For the imputation, I specified an ordered logit model with the other taste variables as predictors of the missing response for each genre the respondent reported not being familiar enough with to produce a response.\footnote{I used the multiple imputation via chained equations routine (\texttt{mi impute chained} command) in \texttt{Stata 15}.} I generated five imputed datasets and assigned respondents the median imputed value for their evaluation of that genre. The distribution of ``not familiar'' responses is highly skewed: $1,511 (67\%)$ of respondents claimed to be familiar with all genres; another $11\%$ reported only one ``not familiar'' response. By the time we get to six ``not familiar'' responses we have reached $95\%$ of respondents. Thus, only 5\% of the analytic sample contains seven or more imputed responses. After dropping respondents who did not provide taste responses on any of the 20 genres ($N = 20$) and dropping respondents who provided the same response across all genres ($N = 21$), which correlational class analysis would treat as its own ``degenerate'' schematic class, the final analytic sample is $N= 2,187$.

\subsection{Correlation Networks for Internal Schema Validation} 
As is now standard, I present the resulting schematic classes as correlation networks \citep{ boutyline2017improving, goldberg2011mapping, epskamp2012qgraph}, with the edge weight set to the estimated correlation for each genre pair for members of that class and the node color set to the assignment of each genre to one of four latent factors based on the binary consumption items (results from which are shown in figure~\ref{fig:typology}). Correlation networks for each taste schema are shown in Figures~\ref{fig:corrnet}a-~\ref{fig:corrnet}e. 

\subsubsection{Expected Influence Score}
To quantify the importance of the evaluation of different genres within each class, I follow the approach recommended by \citet{robinaugh2016identifying} for analyzing node centrality in signed correlation networks and rank each genre in each taste schema by its \emph{expected influence}, which is given by:

\begin{equation}
EI_{gs} = \sum_{j=1}^{N} = r_{gjs}    
\end{equation}

Where $EI_{gs}$ is the expected influence of the $g^{th}$ genre in the $s^{th}$ taste schema and $r_{gjs}$ is the value of the polychoric correlation coefficient between the taste evaluations for genre $g$ and genre $j$ among respondents who share taste schema $s$. This index produces a rank of genres for each schematic class such that those with strong positive correlations with other genres are expected to be the most influential, genres with strong negative correlations are expected to be the least influential, and genres with a mixture of negative and positive correlations are expected to be moderately influential. Figure~\ref{fig:centrality} shows the rank of each genre according to $EI_g$ for each taste schema.\footnote{$EI_g$ is computed using the \texttt{expectedINF} function included the \texttt{R} package \texttt{networktools} \citep{networktools}.}

\section{Results}
\subsection{Number of Taste Schemas and Class Sizes} 
Subjecting the data to CCA, yields a solution (setting weights in the respondent by respondent correlation network below $p<0.01$ to zero) that partitions the sample into five schematic classes.\footnote{To perform correlational class analysis, I used the routines included in the \texttt{R} package \texttt{corclass} \citep{corclass}.} This is similar to the results reported in \citet{boutyline2017improving} applying the same method to General Social Survey data, who found four schematic classes characterizing the American population in 1993. The classes that I obtain using the SSI 2012 data are somewhat uneven in size, with $N = 448 (21\%), 361 (17\%), 700 (32\%), 319 (15\%),$ and 359(16\%) respectively. 
      
\subsection{Taste Schemas in 2012 are Similar to Those Uncovered in 1993 General Social Survey}
The most striking result is that the schematic classes evident in the 2012 data aren very similar to those uncovered by \citet{goldberg2011mapping} and \citet{boutyline2017improving}. This suggests a rather remarkable level of schematic continuity in the American population (at least when it comes to musical genre evaluation). This supports the main implication of the durability position in cultural analysis.

First, one of the schematic classes (Figure~\ref{fig:corrnet}a), which is similar to Boutyline's``anything (but) country,'' is characterized by an opposition of country music to a large clique of positively connected genres (with country featuring a long tie to Gospel/Religious). The key difference between my results and those reported by \citet{boutyline2017improving} is that country does not have numerous negative connections (as in General Social Survey 1993) featuring instead mostly absent or very weak connections. It is clear, however, that members of this schematic class decouple (if not oppose) their evaluations of Country from that of other genres. In addition, just like Boutyline (2017, Figure 10D), I find that Gospel/Religious music is also decoupled from the evaluations of the main cluster in this taste schema, featuring many weak connections (both positive and negative. The main difference is the status of Hip/Hop Rap as a second peripheral genres in this schema, which was not the case in the 1993 data. In fact, as figure~\ref{fig:centrality}a, shows, in these data this schema would be best referred to as Anything (but Country and Rap) as these are the most peripheral genres in the schema. Overall, however, we see a strong continuity pattern in relation to Boutyline's results.  

Second, another of the schematic classes bears a strong resemblance to that which Boutyline (2017, Figure 10) referred to as ``anything (but) heavy metal'' (Figure~\ref{fig:corrnet}b). My results, however, bring out something that was already evident in Boutyline's results: The fact that this schema contrasts the evaluative judgment of both heavy metal {\em and} rap (whose evaluations are strongly connected) against (mostly) everything else. However, as shown in figure~\ref{fig:centrality}b, it is fair to refer to this as a schematic class opposing heavy metal univores against eclectics who prefer other genres, as Metal has the least influence in the correlation network by far. Overall, we once again observe a strong continuity pattern in relation to Boutyline's results.  

Third, just like \citet{boutyline2017improving} and \citet{goldberg2011mapping} I uncover a taste schema that contrasts ``traditional'' (Country, Gospel/Religious, Bluegrass, Folk, Classical, Big Band, Opera) against ``contemporary'' genres (Figure~\ref{fig:corrnet}c). Like Boutyline and Goldberg, I find that Contemporary Rock, Metal, and Rap and Reggae (in addition to Pop, Dance, and Alternative, which were not included in General Social Survey 1993) are squarely on the ``contemporary'' side of the divide. I also find that Blues/R\&B and Latin are in an intermediary (``brokerage'') position between tradition and contemporaneity. The main difference concerns the position of Jazz, which in Boutyline's analysis fell on the side of contemporaneity. In 2012, this genre is more unambiguously on the tradition side. Overall, it is clear that a schematic class following a ``contempo/trad'' logic is a durable feature of the American taste public. 

Fourth, like both Boutyline and Goldberg, my analysis also reveals a schematic class composed of mostly strong and positive, as given by the thick gray edges, and a few null or very weak negative correlations. Like these predecessors, I label this taste schema ``Omnivore-Univore'' although this label is partially misleading, if it is taken to mean that this is the only schematic class capable of producing a version of omnivorousness at the behavioral level; a better assumption to make is that this is the class with the potential for the most extreme forms of omnivorousness in terms of actual consumption choices, a claim we will evaluate below. One thing that was already evident in Boutyline (2017, figure 10) and Goldberg (2011, Figure 3A) but that it is clear in these results, is that the peripheral status of Metal on this schema in relation to the rest of the genres. This is confirmed in Figure~\ref{fig:centrality}d. Which suggests that the prototypical univore in this class is a ``Metalhead'' opposed to persons displaying a favorable attitude to a wide range of musical styles. As before, what is striking here is the similarity of these results to those obtained by Boutyline and Goldberg using data collected about a quarter of a century ago. 

Finally, unlikely Boutyline, but like Goldberg, I uncover a class that seems to make a strong binary partition between most of the meta-genres classified as ``Art,'' ``Folk'' and ``Pop'' and the rest. This class bears a strong resemblance to the one Goldberg labeled ``High-Low'' in General Social Survey 1993 using relationality as a distance criterion, but which Boutyline failed to uncover using the correlation distance in the same data. However, unlike Goldberg, who uncovered a largely binary partition, members of this class seem to divide genres according to a triparte Art/Pop/Folk schema with strong correlations between genres that seem to belong to distinct consumption domains (as shown in~\ref{fig:typology}) (except they do not distinguish between different forms of Pop). These results redeem the existence of a taste schema shared by persons that engage in a more traditional style of classification across meta-genre discourses similar to the traditionalists discussed in \citet{van2001social}. Following this work, I label this schema ``Art/Folk/Pop.'' Taken together, these results suggest substantial levels of schematic continuity in the musical domain, even in the face of large scale transformations. 

\subsection{Internal Validation of Taste Schemas Using Music Listening Choices}
As noted earlier, previous work using schematic class analysis generally lacks a behavioral criterion with which to validate results. Analysts usually have to rely on indirect validation procedures, such as looking at the correlation between taste schemas and demographics \citep{goldberg2011mapping, baldassarri2014neither} and inferring subsequent behavior from what is known about the correlation between demographics and action. The 2012 survey asked respondents to not only evaluate genres but also to list whether they listened to those genres or not. This information provides a source of internal validation for the schematic classes. 

\subsubsection{The Omnivore-Univore Taste Schema Displays the Most Ommnivorousness}
First, we begin by looking at the distribution of the number of consumption choices across taste schemas. This is shown in Figure~\ref{fig:money}. If the omnivore/univore schema lives up to its label, then people who listen to a large number of genres should be over-represented in this class. Figure~\ref{fig:money} shows that this is indeed the case. However, the figure also shows that all taste schemas have some behavioral potential for both omnivorousness and univorousness. In this sense, people who share the same taste schema can end up on either side of the behavioral consumption divide. The key difference will be that different taste schemas will define distinct populations of omnivores and univores who stand opposed in terms of overt choices by virtue of sharing the same logic of taste.

\subsubsection{Centrality in a Taste Schema Predicts the Link Between Omnivorousness and Evaluation}
A more concrete way to validate schematic classes internally is to connect the importance of each genre in the correlation network to the expected consequences its evaluation would have for the expected omnivorousness of people who share that schema. 

We can think of the correlation networks depicted in Figures~\ref{fig:corrnet}a-e as defining a ``spreading activation'' network for each schematic class \citep{collins1975spreading}. The concept of spreading activation is simple. Imagine picking a genre at random and ``activating it'' by presuming that somebody likes that genre. If this genre is also linked to a lot of other genres via positive correlation edges, then the evaluative activation should spread to those other genres, such that liking the focal genre will result in people who share that schema also liking the other musical styles to which it is positively linked. Liking lots of genres should result in an increased likelihood of actually listening to a lot of musical styles under normal conditions. 

Now imagine, the same process but this time we ``activate'' (by producing a positive evaluation) a genre that is peripheral in the correlation network of that schema (like Country or Hip Hop in Figure~\ref{fig:corrnet}a). Because such genres feature mostly null or negative links with other musical styles, then we should expect that the evaluation of those genres should be decoupled (or in the case of negative correlations lead to a devaluation) of the other genres in the schema. In this case, there is no spreading activation and no empirical expectation that evaluations of these genres should be linked to behavioral omnivorousness. 

If this imagery of the way taste schemas link to behavior is on the right track, we should expect that behavioral omnivorousness should be more strongly linked to the positive evaluation of central genres in each taste schema. In the same way, the number of genres actually listened to should have a lower correlation with the evaluation of genres that are peripheral in the correlation network for that schema. This is a cognitively plausible, and empirically verifiable, way of conceptualizing the link between taste schemas and consumption behavior. 

Results of an analysis meant to evaluate this conjecture are shown in Figure~\ref{fig:validation1}. In the figure the y-axis gives the Spearman rank correlation between the number of genres that people who share that schema listen to (as given in figure~\ref{fig:money}) and the ordinal taste evaluation item for that genre (higher values coded as liking that genre), rescaled to the variance for that taste schema. The X-axis is the expected influence ($EI_{gs}$) score of the genre in the corresponding taste correlation network (as given in figure~\ref{fig:centrality}). 

As expected, we see a strong positive correlation between these two quantities within all five taste schemas. The upper right of each scatterplot shows that there is a stronger than expected correlation between the number of genres listened to and the evaluation of genres that are central to the schema. The lower-left quadrant of each scatterplot shows that there is a lower than expected correlation between behavioral omnivorousness and the evaluations of genres that are peripheral to each taste schema. This constitutes strong evidence for the internal validity of the correlation networks extracted from each taste schema and the spreading activation interpretation of the link between the resulting schemas and cultural choice behavior. 

\subsubsection{Centrality in a Taste Schema Predicts Likelihood Omnivores will Consume the Genre}
As Figure~\ref{fig:money} shows, some schematic classes have greater potential for omnivorousness. However, all classes display some level of omnivorousness. The question is what is the difference in the omnivore repertoire of the different schematic classes? If the correlation network of each class encodes information relevant to actual cultural choice behavior, and if evaluations and number of actual choices are most closely linked to the most central genres in this network, we should expect this genres to also be represented in the repertoire of omnivores in each taste class, resulting in distinct behavioral expressions of this taste pattern. 

Figure~\ref{fig:validation2} shows that due to the substantive correlation between omnivorousness and evaluation, centrality in the taste schema correlation network can be used to derive the different ``forms of omnivorousness'' that we should expect each taste schema to yield. In the figure, the y-axis corresponds to the difference between the proportion of omnivore respondents in that class (indexed as those respondents who listen to more than eight genres, which corresponds to the $90^{th}$ percentile of the sample distribution) who listen to that genre and the overall popularity of the genre across all respondents. The length of the dropline from left to right thus gives us a sense of the extent to which certain genres are likely to be over-represented in the taste repertoire of behavioral omnivores in each schematic class. As can be readily verified, with some exceptions, genres over-represented in the taste repertoire of omnivores are also most central in the correlation network for that class. The scatter plot shown in Figure~\ref{fig:validation3}a, shows that this is generally the case, as the two quantities exhibit a moderately strong correlation across taste schemas ($r=0.49, p<0.01$). 

Figure~\ref{fig:validation3}b, confirms something that is evident in Figure~\ref{fig:validation2}, and that is that the correlation between genre centrality in the taste schema and its over-representation in the schema-specific omnivore repertoire is stronger for those genres that feature more inequality in their relative chances of being engaged by omnivores across schematic classes. The Figure splits the top scatter plot by the standard deviation of the over-representation score across the five classes. The left-most panel shows genres in the lowest standard deviation quartile and thus with relatively equitable probabilities of being present in the omnivore repertoire of different taste schemas. The right-most panel shows genres in the top quartile, and thus with the most heterogeneous representation in the omnivore repertoire of different taste schemas. As the figure shows, the centrality of the genre in the last group is a strong predictor of their showing up in the listening repertoire of omnivores in each class. 

\section{Discussion and Conclusion}
This paper began with the debate over the malleability versus the durability of schemas in cultural analysis. To shed empirical light on this issue, I revisited the relatively well-studied issue of heterogeneity in evaluative schemas for musical genres in the American population \citep{goldberg2011mapping}. Using recently collected survey data comparable to the General Social Survey 1993 culture module, I find that taste schemas of Americans in 2012, bear a strong resemblance from those uncovered in previous work using comparable techniques (most closely \citet{boutyline2017improving}). I interpret these results as providing strong evidence that, at least at the population level, there is strong continuity and durability in the way Americans organize their understandings of an important cultural domain. These findings deepen and extend recent arguments arguing against extreme versions of the schematic malleability thesis, showing that we find relatively strong levels of population-level durability in a variety of cultural patterns \citep{vaisey2016cultural}.

This paper also contributes and extends previous work in survey-based schematic class analysis by attempting to validate the results of the resulting schema partitions using criteria internal to the domain in question (in this case, patterns of musical listening choices). I showed that we can leverage techniques developed for the analysis of correlation networks to go beyond eyeballing \citep{epskamp2012qgraph,networktools}. In particular, I showed that a simple measure of vertex centrality referred to as ``expected influence'' \citep{robinaugh2016identifying}, does a good job of capturing the genres most central in each schema. The validation results reveal two sets of notable findings. 

I found that the taste schema featuring the lowest number of negative links among evaluative items (labeled Omnivore-Univore in previous work) is indeed the most omnivorous in terms of their actual likelihood to listen to a wide variety of musical styles. This supports the idea that at least in some cases, behavioral predictions can be made from ``eyeballing'' the correlational network corresponding to a given schema. However, beyond these broad predictions, it is perilous to conflate the evaluative structure of a schema with its behavioral potential. For instance, I find that all taste schemas have potential for a form of omnivorousness, but this cannot be ascertained from a global assessment of the correlation structure. Instead, the relative importance of genres within the schema must be linked to behavioral outcomes in question. 

Following this approach, I showed that influential genres in the correlation network are also those for which there is a strong correlation between omnivorousness (in terms of actual choice behavior) and positive evaluation. I argued that this pattern of results can be accounted for if we read the schema correlation as a possible conduit for ``spreading activation'' of evaluations. This allows us to pick the genres the evaluation of which will be strongly coupled to the evaluation of other genres in the schema, from those whose evaluation will be loosely coupled. 

Third, I showed that genres central in the correlation network of the taste schema, are also the ones that will be most likely to be found in the consumption portfolio of omnivores who share that schema. This is most likely to be the case for those genres that have the least equitable distribution among omnivores across schemas. A possible explanation for this result is that cultural objects whose evaluations and engagement behavior are most driven by schema-driven evaluations (and which therefore end up having an unequal distribution among people with different taste schemas) are also the ones whose differential engagement is most predictable from the structure of evaluations in the schema. These objects are different from those whose consumption may be driven by external factors and which thus become ``prescribed'' no matter what, driving down variance in engagement among active members across schematic classes and attenuating the link between position in the correlation network and actual behavior. This can be a fruitful avenue for further research in schematic class analysis. 


\newpage
\section*{References}
\bibliographystyle{elsarticle-harv}
\bibliography{main.bib}
\newpage

%Factor analysis of consumption items
\begin{figure}[t!]
\centering\includegraphics[width=1.0\linewidth]{genre-behav-typology}
\caption{\scriptsize Behavioral Meta-Genre Typology Obtained from Factor Analysis of Musical Listening Choices. Circular Nodes are Latent Factors and Square Nodes are Genre Labels. Green Edges Indicating Positive Loadings on the Factor, Red Edges Indicate Negative Loadings (or cross-factor correlations). Legend: Cnt = Country, Blg = Bluegrass, Flk = Folk, Gsp = Gospel/Religious, Ltn = Latin/Spanish/Salsa, Rap = Hip Hop/Rap, Bls = Blues and R\&B, Rgg = Reggae, Dnc = Dance/Electronic, Cls = Classical, Opr = Opera, Jzz Jazz, Brd = Broadway/Musicals, Esy = Easy Listening, Bgb = Big Band, ClR = Classic Rock/Oldies, Pop = Pop/Top 40, CnR = Contemporary Rock, Alt = Alternative Rock, Mtl = Hard Rock/Heavy Metal.}
\label{fig:typology}
\end{figure}

%comparing standard and LASSO regularized correlation networks (node-link diagrams)
\newpage
\begin{figure}[ht!]
  \centering
  \captionsetup{justification=centering}
    \subfloat[][\scriptsize Anything (but Country) \\ N = 478]{\includegraphics[width=0.33\textwidth]{taste-corrnet-class1}} 
    \subfloat[][\scriptsize Anything (but Heavy Metal) \\ N = 365 ]{\includegraphics[width=0.33\textwidth]{taste-corrnet-class2}} 
    \subfloat[][\scriptsize Contemporary/Traditional \\ N = 619]{\includegraphics[width=0.33\textwidth]{taste-corrnet-class3}} \\
    \subfloat[][\scriptsize Omnivore/Univore \\ N = 376]{\includegraphics[width=0.33\textwidth]{taste-corrnet-class4}} 
    \subfloat[][\scriptsize Contempo/Folksy/Easy \\ N = 363]{\includegraphics[width=0.33\textwidth]{taste-corrnet-class5}} \\
    \subfloat[][\scriptsize Standardized Expected Influence Scores for Each Genre]{\includegraphics[width = 1.0\textwidth]{taste-centrality-class}} 
  \captionsetup{justification=justified}
  \caption{\scriptsize \newline \textbf{Top (a-e):} Correlation Networks for each taste schema. Node color corresponds to the assignment of the genre to one of four latent factors based on independent consumption items (see Figure~\ref{fig:typology}). Edge color indicates a positive (gray) or negative (red) correlation between genre evaluations within each class. Edge thickness is proportional to the absolute value of the polychoric correlation. Correlations that do not meet a statistical significance threshold ($p<0.01$) are set to zero. Significant correlations below an effect size threshold ($r < 0.05$) are omitted from the plot to increase readability. \newline \textbf{Bottom (f):} X-axis is the (standardized at taste schema class mean) expected influence score $EI_{gs}$ for that genre. For each taste schema, genres placed toward the top have the most expected influence within that schema. Genres located toward the bottom are most peripheral.} 
  \label{fig:corrnet}
 \end{figure}
    
 %N. of Genres by Schematic Class
 \newpage
\begin{figure}[ht!]
  \centering
    \subfloat[][]{\includegraphics[width = 0.8\textwidth]{numlis-by-schematic-class.png}} \\
    \subfloat[][]{\includegraphics[width = 1.0\textwidth]{schema-validation1.png}} 
    \caption{\newline\textbf{Top:} X-Axis is the number of genres respondents report listening to. Y-axis is the proportion of respondents in the sample who fall into that category. \newline\textbf{Bottom:} Y-axis is the (rescaled within taste schema) Spearman rank correlation between number of genres respondents listen to and their taste evaluation of that genre (higher values indicating the respondent likes the genre). X-axis is the (standardized at taste schema class mean) expected influence ($EI_{gs}$) of that genre in the correlation network for that schema (as given in figure~\ref{fig:corrnet}f). Within each taste schema, the positive correlation between omnivorousness and evaluation increases as the centrality of that genre in the schema goes up. Spearman rank correlation coefficient per panel: Anything but Country: $r=0.62, p<0.01$; Anything but Heavy Metal: $r=0.70, p<0.01$; Contempo-Trad: $r=0.68, p<0.01)$; Omnivore-Univore $r=0.90, p <0.01 $; Contempo/Easy/Folksy $r=0.61, p<0.01$.}
  \label{fig:money}
 \end{figure}
 
 \newpage
 \begin{figure}[ht]
 \includegraphics[width = 1.0\textwidth]{schema-validation2.png}
 \caption{}
 \label{fig:validation2}
 \end{figure}
 
 \newpage
 \begin{figure}[ht]
 \includegraphics[width = 1.0\textwidth]{schema-validation5.png}
 \caption{\newline\textbf{Left:} Y-axis is the extent to which the genre is over-represented in the omnivore consumption profile of individuals in each class (as given in Figure~\ref{fig:validation2}). X-axis is the centrality of the genre in the correlation network (as given in Figure~\ref{fig:corrnet}f). Across each taste schema, the likelihood of a genre being engaged by omnivores is correlated with the centrality of the genre in the taste schema (Pearson's $r=0.49, p<0.01$).\newline \textbf{Right:} Y-axis is the extent to which the genre is over-represented in the omnivore consumption profile of individuals in each class (as given in Figure~\ref{fig:validation2}). X-axis is the centrality of the genre in the correlation network (as given in Figure~\ref{fig:corrnet}f). Panels are arranged by the standard deviation of omnivore genre over-representation across the five taste schemas. Upper-left panel is the bottom quartile (most equitably distributed across schemas) and the bottom-right panel is the top quartile (least equitably distributed across schemas). The correlation is stronger for genres that have unequal representation in the omnivore consumption profile of each taste schema. Pearson correlation coefficient per panel: $r=0.12, p=0.55$; 2: $r= 0.40, p<0.05$; 3: $r=0.47, p<0.05)$; 4 $r=0.64, p\leq0.01 $.}
 \label{fig:validation3}
 \end{figure}
 
%% Authors are advised to submit their bibtex database files. They are
%% requested to list a bibtex style file in the manuscript if they do
%% not want to use model1-num-names.bst.

%% References without bibTeX database:

% \begin{thebibliography}{00}

%% \bibitem must have the following form:
%%   \bibitem{key}...
%%

% \bibitem{}

% \end{thebibliography

\end{document}

%%
%% End of file `elsarticle-template-1-num.tex'.